% swift-idiom-patterns.tex — the patterns Swift developers actually write that are not in GoF:
% type erasure (hand box vs any P vs some P), phantom types, newtype / tagged IDs, protocol
% witnesses (struct of closures, as in pointfreeco/swift-dependencies), result builders vs fluent
% API, closures as Strategy, enums as state machines, default implementations + the subclass
% dispatch trap, generic constraints as compile-time strategy.
% Sources (checked 2026-10-06):
%   swift-evolution 0335 (any; ExistentialAny upcoming feature) · 0309 (existentials for protocols
%     with associated types / Self, Swift 5.7 — these require `any`) · 0346 (primary associated
%     types, 5.7) · 0352 (implicit opening of existentials, 5.7) · 0289 (result builders, 5.4)
%   github.com/swiftlang/swift docs/ABI/TypeLayout.rst (existential container: 3-word inline
%     value buffer + type metadata + one witness table per protocol)
%   github.com/pointfreeco/swift-tagged (Tagged<Tag, RawValue>) · github.com/pointfreeco/
%     swift-dependencies DesigningDependencies.md (struct-of-closures clients)
%   developer.apple.com/documentation/swiftui/anyview (type erasure loses structural identity)
% @labels: area=design kind=pattern platform=apple topic=language,patterns discussion=103
% @tags: type-erasure, existentials, witness-table, generic-specialization, primary-associated-types, anyview, phantom-types, tagged-ids, type-state, protocol-witnesses, result-builders, enum-state-machine
\documentclass[8pt]{extarticle}
\usepackage{printup-sheet}

\lstdefinelanguage{SwiftSheet}{
  morekeywords={protocol,class,final,struct,enum,func,var,let,weak,init,case,switch,
    if,else,return,guard,self,Self,nil,try,await,async,throws,private,static,override,
    true,false,Void,String,Bool,Int,Data,Date,UUID,Duration,extension,where,some,any,
    associatedtype,mutating,for,in,do,catch,default,Sendable},
  sensitive=true, morecomment=[l]{//}, morecomment=[s]{/*}{*/}, morestring=[b]",
  literate={->}{{\hbox{-}\hbox{>}}}2 {==}{{\hbox{=}\hbox{=}}}2 {??}{{\hbox{?}\hbox{?}}}2
           {!=}{{\hbox{!}\hbox{=}}}2 {<<}{{\hbox{<}\hbox{<}}}2}

\newcommand\pat[3][sheetBlue]{\par\noindent\fcolorbox{#1}{#1!5}{\parbox{\dimexpr\linewidth-2\fboxsep-2\fboxrule\relax}{\footnotesize\raggedright{\bfseries\color{#1}#2}\enspace #3}}\par\vspace{1pt}}

\tikzset{
  w/.style={cell, minimum width=24mm, minimum height=4mm, font=\ttfamily\tiny},
  sb/.style={box, font=\scriptsize, inner sep=1.5pt, minimum height=4.2mm},
  lbl/.style={font=\tiny, text=black!80, inner sep=1pt, align=center},
  ttl/.style={font=\bfseries\small, text=sheetBlue},
}

\begin{document}

\sheettitle{Swift idiom patterns — beyond GoF}{design · folio}

\oneliner{Swift turns runtime patterns into \textbf{compile-time} ones: a \textbf{phantom type}/\textbf{tagged
ID} makes a wrong argument a compile error, an \textbf{enum} makes an illegal state unrepresentable, a
\textbf{generic constraint} is a Strategy chosen at build time — and \textbf{type erasure} is the price
of going back to runtime. Order of preference: \textbf{\texttt{some}/generics $>$ \texttt{any P<X>} $>$
a hand-written \texttt{AnyX}}.}

\vspace{3pt}
\noindent\begin{tikzpicture}[sheet]
  \foreach \x in {4.3,8.55,12.6} \draw[sheetGrey!40] (\x,3.05) -- (\x,-0.2);
  % ── existential
  \node[ttl] at (2.05,2.9) {\texttt{let s: any Shape}};
  \node[w, fill=sheetBlue!6] (b1) at (1.4,2.3) {buffer: 3 words inline};
  \node[w, fill=sheetBlue!6] (b2) at (1.4,1.9) {type metadata $\to$ Circle};
  \node[w, fill=sheetOrange!12, draw=sheetOrange] (b3) at (1.4,1.5) {witness table (PWT)};
  \node[sb, draw=sheetRed, fill=sheetRed!6, font=\tiny] (hp) at (3.55,2.3) {heap box\\if $>$ 3 words};
  \draw[->, draw=sheetRed] (b1) -- (hp);
  \node[sb, fill=sheetOrange!10, draw=sheetOrange, font=\tiny] (wa) at (3.55,1.5) {\texttt{area} $\to$\\Circle.area};
  \draw[hot] (b3) -- (wa);
  \node[lbl] at (2.05,0.75) {runtime box · any conformer\\call = load PWT, jump · \textbf{no specialisation}};
  \node[lbl, text=sheetBrown] at (2.05,0.1) {\texttt{[any Shape]} can mix types};
  % ── unspecialised generic
  \node[ttl] at (6.4,2.9) {\texttt{func f<T: Shape>(\_ x: T)}};
  \node[sb] (cal) at (5.35,2.2) {caller: \texttt{f(circle)}};
  \node[sb, minimum width=24mm, fill=black!4, draw=sheetGrey] (bod) at (6.4,1.1) {ONE compiled body};
  \draw[flow] (cal) -- (bod);
  \node[lbl, anchor=west, align=left] at (6.5,2.2) {hidden args:\\\textbf{T's metadata}\\+ \textbf{T: Shape PWT}};
  \node[lbl] at (6.4,0.55) {no box — but still witness calls};
  \node[lbl, text=sheetBrown] at (6.4,0.1) {one T per call: \texttt{[T]} can't mix};
  % ── specialised
  \node[ttl] at (10.55,2.9) {specialised by the optimiser};
  \node[sb, draw=sheetGreen, fill=sheetGreen!10] (s1) at (9.6,2.1) {\texttt{f<Circle>}};
  \node[sb, draw=sheetGreen, fill=sheetGreen!10] (s2) at (11.5,2.1) {\texttt{f<Square>}};
  \node[sb, fill=black!4, draw=sheetGrey, font=\tiny] (g) at (10.55,1.25) {generic \texttt{f<T>}};
  \draw[flow] (g) -- (s1); \draw[flow] (g) -- (s2);
  \node[lbl] at (10.55,0.45) {cloned per concrete type $\to$ \textbf{direct},\\inlinable calls (same module, WMO,\\\texttt{@inlinable}); \texttt{some P} gets this too};
  % ── hand box
  \node[ttl] at (14.6,2.9) {hand-written \texttt{AnyShape}};
  \node[sb, minimum width=24mm] (ab) at (14.6,2.25) {\texttt{struct AnyShape}\\\texttt{\_area: () $\to$ Double}};
  \node[sb, fill=sheetOrange!10, draw=sheetOrange, font=\tiny] (ctx) at (14.6,1.25) {closure context\\(heap) holds \textbf{Circle}};
  \draw[hot] (ab) -- (ctx);
  \node[lbl] at (14.6,0.55) {concrete type, erased by closures;\\\texttt{AnyPublisher}, \texttt{AnySequence} work so};
  \node[lbl, text=sheetBrown] at (14.6,0.1) {costs an allocation + indirect call};
\end{tikzpicture}

\begin{multicols}{2}
\raggedright

\pat{Type erasure — three tools}{\textbf{Hand box}: store the conformer's methods as closures
$\to$ a \emph{concrete} type you can extend and conform (\texttt{AnyPublisher},
\texttt{AnySequence}). \textbf{\texttt{any P<X>}} (primary associated types, 5.7): a built-in box.
\textbf{\texttt{some P}}: \emph{not} erasure — hidden from the reader, known to the compiler.
\texttt{eraseToAnyPublisher()} hides \texttt{Publishers.Map<…>}; \texttt{AnyView} also erases
\textbf{structural identity} (worse diffing, lost \texttt{@State}); \texttt{AnyHashable} = mixed
keys in one \texttt{Set}. An \texttt{any P} passed to \texttt{<T: P>} is \textbf{opened} (SE-0352).}
\begin{lstlisting}[language=SwiftSheet]
protocol Loader<Output> { associatedtype Output
  func load() async throws -> Output }
struct AnyLoader<Output>: Loader {         // hand-written box
  private let _load: () async throws -> Output
  init<L: Loader>(_ base: L) where L.Output == Output {
    _load = base.load }                    // captures base
  func load() async throws -> Output { try await _load() } }
let b: any Loader<User> = RemoteUserLoader() // built-in (5.7)
\end{lstlisting}

\pat{Phantom types · tagged IDs · newtype}{A \textbf{phantom} parameter is never stored; it only
stops mixing values with the same representation: \textbf{tagged IDs} (pointfree's
\texttt{Tagged<Tag, Raw>}), \textbf{type-state} (a \texttt{Request<Draft>} has no \texttt{send()}).
\textbf{Newtype}: \texttt{struct Email} with a failable \texttt{init} validates \emph{once} —
``parse, don't validate''. Zero runtime cost.}
\begin{lstlisting}[language=SwiftSheet]
struct ID<Entity>: Hashable, Sendable { let raw: UUID }
struct User { let id: ID<User> }
struct Order { let id: ID<Order>; let buyer: ID<User> }
func user(_ id: ID<User>) -> User? { nil }
// user(order.id)   error: ID<Order> is not ID<User>
enum Draft {}; enum Signed {}              // uninhabited tags
struct Request<State> { var body: Data }
extension Request where State == Draft {
  func signed() -> Request<Signed> { .init(body: body) } }
extension Request where State == Signed { func send() async {} }
\end{lstlisting}

\pat{Protocol witnesses — a struct of closures}{One struct whose \emph{properties are closures}
replaces protocol + conformers; each ``conformance'' is a static value (\texttt{.live},
\texttt{.mock}). \textbf{+} override \emph{one} endpoint in a test, derive variants by transforming
values, no \texttt{any}/associated-type friction. \textbf{$-$} no default methods, no \texttt{Self};
closures must be \texttt{@Sendable} to share. pointfree's \textbf{swift-dependencies} clients are
built this way.}
\begin{lstlisting}[language=SwiftSheet]
struct TimeClient: Sendable {
  var now: @Sendable () -> Date
  var sleep: @Sendable (Duration) async throws -> Void }
extension TimeClient { static let live = TimeClient(
  now: { Date() }, sleep: { try await Task.sleep(for: $0) }) }
var t = TimeClient.live; t.now = { .distantPast } // one endpoint
\end{lstlisting}

\pat{Builder: result builder vs fluent API}{\texttt{@resultBuilder} (SE-0289): the compiler rewrites
the closure's statements into \texttt{buildExpression}/\texttt{buildBlock}/\texttt{buildOptional}/
\texttt{buildEither}/\texttt{buildArray} calls — a compile-checked DSL (\texttt{@ViewBuilder}).
\textbf{Fluent}: each call returns a modified \emph{copy} of a struct —
\texttt{Req().header(k, v).timeout(5)}; no \texttt{build()}, never half-built.}
\begin{lstlisting}[language=SwiftSheet]
@resultBuilder enum PathBuilder {
  static func buildExpression(_ s: String) -> [String] { [s] }
  static func buildBlock(_ c: [String]...) -> [String] {
    c.flatMap { $0 } }
  static func buildOptional(_ c: [String]?) -> [String] { c ?? [] } }
func path(@PathBuilder _ make: () -> [String]) -> String {
  make().joined(separator: "/") }
let p = path { "users"; if isAdmin { "admin" } } // users/admin
\end{lstlisting}

\columnbreak

\pat{Closures as Strategy · enums as state machines}{A one-method Strategy is a function:
\texttt{sorted(by:)}, a \texttt{price} closure. Use a protocol once there are \emph{several} related
operations, state, or it needs a name. \textbf{State machine}: each case carries only the data valid
in that state (no \texttt{isLoading} + \texttt{error?} + \texttt{items} combinations); transitions
= one \texttt{switch} on \texttt{(state, event)}.}
\begin{lstlisting}[language=SwiftSheet]
enum Load { case idle, loading, loaded([Item]), failed(String) }
enum Event { case start, success([Item]), failure(String) }
extension Load { mutating func send(_ e: Event) {
  switch (self, e) {
  case (.idle, .start), (.failed, .start): self = .loading
  case (.loading, .success(let v)):         self = .loaded(v)
  case (.loading, .failure(let m)):         self = .failed(m)
  default: break } } }   // illegal transition: assert/log here
\end{lstlisting}

\pat{Default implementations — the subclass trap}{Extensions give defaults and mixins. If a
\textbf{superclass} conforms \emph{using} the default, the witness is bound for \texttt{Base};
\texttt{Sub}'s method (no \texttt{override} needed!) is never reached through the protocol.
\textbf{Fix}: implement it in \texttt{Base}; \texttt{Sub} must then write \texttt{override}.}
\begin{lstlisting}[language=SwiftSheet]
protocol Greeter { func hi() -> String }
extension Greeter { func hi() -> String { "default" } }
class Base: Greeter {}                 // witness = the extension
class Sub: Base { func hi() -> String { "sub" } } // no override
let g: any Greeter = Sub()
g.hi()   // "default"  (but Sub().hi() is "sub")
\end{lstlisting}

\pat{Generic constraint = compile-time Strategy}{The strategy is a \emph{type} fixed at the use
site: specialised and inlined, no stored existential, no runtime swap.}
\begin{lstlisting}[language=SwiftSheet]
protocol Backoff { static func delay(_ n: Int) -> Duration }
enum Exponential: Backoff { static func delay(_ n: Int)
  -> Duration { .milliseconds(100 << n) } }
struct Retrier<B: Backoff> {
  func wait(attempt n: Int) async throws {
    try await Task.sleep(for: B.delay(n)) } } // static call
let r = Retrier<Exponential>()         // fixed at build time
\end{lstlisting}

\section{Traps}
\begin{itemize}
  \trap{\texttt{AnyView} everywhere: lost identity $\to$ worse diffing, reset \texttt{@State}.}
  \trap{``\texttt{some} is type erasure'' — no: the compiler still knows the type.}
  \trap{Bare \texttt{P} still means \texttt{any P}; \texttt{any} is required for associated-type/\texttt{Self}
        protocols, everywhere with \texttt{ExistentialAny}.}
\end{itemize}

\section{Remember}
\textbf{Make wrong code not compile:} phantom $\to$ wrong ID · enum $\to$ wrong state · newtype
$\to$ unvalidated input. \textbf{\texttt{some} $>$ \texttt{any} $>$ \texttt{AnyX}.}


\end{multicols}

\noindent{\footnotesize\color{sheetGrey}\textit{Related:} protocols-generics (dispatch) ·
patterns-in-apple-frameworks · dependency-injection-deep · gof-behavioural · gof-creational-structural
(Builder) · functional-design}

\end{document}
